% 对应题号 M1-73；图与本题条件、推导、结论逐项对应。
% 图契约：印刷宽16cm；最小字号8.4pt；几何/场或曲线在左，计算关系在右。
\documentclass[tikz,border=7pt]{standalone}
\usepackage{fontspec,xeCJK,amsmath,amssymb,graphicx}
\IfFontExistsTF{Times New Roman}{\setmainfont{Times New Roman}}{\setmainfont{TeX Gyre Termes}}
\IfFontExistsTF{Heiti SC}{\setCJKmainfont[AutoFakeBold=1.5]{Heiti SC}}{\setCJKmainfont[AutoFakeBold=1.5]{FandolHei-Regular}}
\usetikzlibrary{arrows.meta,bending,calc,patterns,positioning}
\definecolor{ink}{HTML}{243745}
\definecolor{blue}{HTML}{275DAB}
\definecolor{green}{HTML}{177C63}
\definecolor{red}{HTML}{BD4935}
\definecolor{orange}{HTML}{AA6C16}
\definecolor{purple}{HTML}{715393}
\tikzset{every node/.style={font=\fontsize{9.3}{12.5}\selectfont,text=ink},
 axis/.style={-{Stealth[length=4pt]},line width=.65pt,draw=ink!65},
 vec/.style={-{Stealth[length=5pt]},line width=1pt,draw=blue},
 short/.style={-{Stealth[length=3pt]},line width=.8pt},
 guide/.style={dashed,draw=ink!40,line width=.5pt},
 card/.style={rounded corners=2pt,draw=blue!65,fill=blue!4,align=center,text width=6.6cm,inner xsep=9pt,inner ysep=7pt},
 title/.style={font=\fontsize{10.5}{14}\selectfont\bfseries,text width=15cm,align=center},
 label/.style={font=\fontsize{8.4}{11}\selectfont,align=center}}
\begin{document}
\begin{tikzpicture}[x=1cm,y=1cm]
\path[use as bounding box] (-8,-3.4) rectangle(8,3.4);
\node[title] at (0,3) {M1-73\quad 径向场的体源与点源分别记账};

\begin{scope}[shift={(-4.1,0)}]
\draw[ink!45,fill=blue!3] (0,0) circle(1.65);\fill[red] (0,0) circle(2pt);
\foreach \a in {0,45,...,315}{\draw[vec,blue] (\a:.65)--(\a:1.6);\draw[short,orange] (\a:2.15)--(\a:1.8);}
\draw[guide] (0,0)--(45:1.65) node[midway,above left] {$R$};
\node[label,blue] at (0,2.45) {$\beta\mathbf e_r/r^2$：原点集中源};
\node[label,orange] at (0,-2.45) {$-\alpha r\mathbf e_r$：均匀体源项};
\end{scope}
\node[card,draw=blue!65,fill=blue!4] at (3.95,1.65) {$\displaystyle E_r=-\alpha r+\frac\beta{r^2}$\\$\Phi(R)=4\pi(-\alpha R^3+\beta)$};
\node[card,draw=green!65,fill=green!4] at (3.95,0) {$\nabla\cdot\mathbf E=-3\alpha$（$r>0$）\\$\nabla\times\mathbf E=0$};
\node[card,draw=orange!65,fill=orange!4] at (3.95,-1.65) {$\nabla\cdot\mathbf E=-3\alpha+4\pi\beta\delta^{(3)}(\mathbf r)$\\箭头分别示意 $\alpha,\beta>0$ 的两项。};
\end{tikzpicture}\end{document}
